% foundations/part_intro.tex

\chapter*{Mathematical Foundations}
\phantomsection
\addcontentsline{toc}{chapter}{Mathematical Foundations}

The curvature edit has carried the argument to a precise threshold. The thesis
must now describe what remains the same, what changes, which quantities follow,
and under which context an evaluation is valid. Natural-language distinctions
are necessary, but they are not sufficient for derivation, comparison, or
implementation.

This part therefore gives selected engineering dependencies a mathematical
form. It begins with information modelling and conceptual structure because a
formula is only useful when the roles of its inputs and outputs are clear. It
then states the governing axioms and notation, develops state and operator
relations, and explains how intrinsic construction acquires measurable spatial
meaning through metric realization.

Mathematics does not decide the identity of an engineering object, admit
knowledge, or exercise engineering authority. It can make dependencies explicit,
derive states, propagate a curvature change, compare an observation with a
target, and evaluate a constraint. The distinctions established in Part I must
survive every such operation.

The orienting question is:

\begin{quote}
Which formal dependencies make consequences computable without confusing the
alignment, its states, its representations, the evidence about it, and the
decision that may follow?
\end{quote}

The first chapter crosses this threshold through the representations already
encountered in practice.
